\scanpage{311}
\mspage{10}
% archive sheet number: 312
fournit \deleted{un}

\uline{Prop. 2} Avec les notations \inserted{et hyp.} de prop. 1, avec $A$ et $B$ additives \deleted{karoubiennes}

\[
\begin{aligned}
\text{a)}\quad&\underline{\operatorname{Hom}}_{\substack{\text{comm}\\\text{à }\varinjlim}}(\mathrm{Ss},A)\xrightarrow{\sim} A^{.}\\
\text{b)}\quad&\underline{\operatorname{Hom}}_{\substack{\text{comm}\\\text{à }\varprojlim}}((\mathrm{Ss})^o,B)\xrightarrow{\sim} B_{.}
\end{aligned}
\]

Nous allons \deleted{\unsure{préciser}}\inserted{expliciter} cette correspondance du point de vue de l'algèbre universelle. Prenons d'abord le cas de b), avec $B=(\mathrm{Ab})$ pour simplifier. \deleted{\unsure{résultat|résulte}}\inserted{Donc il faut} préciser la relation entre
\[
F:(\mathrm{Ss})^o\xrightarrow[\substack{\text{commute}\\\text{aux }\varprojlim}]{}\mathrm{Ab} \qquad\text{et}\quad K_{.}=(K_n,d_n)\in \mathrm{Ab}_{.}
\]
qui se déterminent mutuellement. On désigne comme de juste par $K_*$ l'objet de $\mathrm{Ab}_*$ défini par $K_{.}$. On a donc
\[
F(X_*)\simeq\operatorname{Hom}_{\widehat{\mathbb{\Delta}}}(X_*,K_*)\simeq\operatorname{Hom}_{\widehat{\mathbb{\Delta}}_{ab}}(\mathbb{Z}(X_*),K_*)
\]
où $\deleted{\unsure{\{|(}}\mathbb{Z}(X_*)\deleted{\unsure{\}|)}}\deleted{\unsure{{}_*|*}}$ désigne le ``faisceau abélien libre engendré par $X_*$'' en l'occurrence
\[
\mathbb{Z}(X_*)=(n\longmapsto\mathbb{Z}(X_n)=\mathbb{Z}^{(X_n)})
\]
Ce groupe \inserted{ab} $1/2$ simplicial sera aussi noté $\mathbb{C}_*(X_*)$, et $\mathbb{C}_{.}(X_*)$ le complexe de chaines associé \inserted{par DP} : \deleted{On aura donc}
\[
\left\{\begin{aligned}
\mathbb{C}_*(X_*)&=\mathbb{Z}(X_*)=(n\longmapsto\mathbb{Z}(X_n)=\mathbb{Z}^{(X_n)}\\
\mathbb{C}_{.}(X_*)&=\mathrm{ND}_{.}(\mathbb{C}_*(X_*))
\end{aligned}\right.
\]
Donc par Dold-Puppe on aura
\[
F(X_*)\simeq\operatorname{Hom}_{\substack{\widehat{\mathbb{\Delta}}_{ab}\\\shortparallel\\\mathrm{Ab}_*}}(\mathbb{C}_*(X_*),K_*)\underset{\text{par D.P.}}{\simeq}\operatorname{Hom}_{\mathrm{Ab}_{.}}(\mathbb{C}_{.}(X_*),K_{.})
\]
Donc le foncteur $F_{K_{.}}=F\quad \mathrm{Ss}^o\longrightarrow \mathrm{Ab}$
\[
X_*\longmapsto F_{K_{.}}(X_*)=F(X_*)
\]
associé à $K_{.}\in \mathrm{Ab}_{.}$ est donné par
\[
\boxed{F_{\unsure{K_{.}|K_*}}(X_*)=\operatorname{Hom}_{.}(\mathbb{C}_{.}(X_*),K_{.})}
\]
On en tire deux conclusions :

1°) Les valeurs sur $X_*$ des foncteurs envisagés $F$ ($\mathrm{Ss}^o\longrightarrow \mathrm{Ab}$, commutant aux $\varprojlim$) ne dépendent que de $\mathbb{C}_*(X_*)\in \mathrm{Ab}_*$ \inserted{(``abélianisé'' de $X_*$)} --- ce qui revient au $=$, de $\mathbb{C}_{.}(X_*)\in \mathrm{Ab}_{.}$ ; \deleted{\unsure{donc le|donc ledit}} complexe de chaines $\mathbb{C}_{.}(X_*)_*$ peut être considéré comme ``le'' invariant additif universel'' de $X_*$. Quand on ``oublie $X_*$ au profit de $\mathbb{C}_{.}(X_*)^{\unsure{*}}=C_{.}\in \mathrm{Ab}_{.}$'', les foncteurs envisagés (en tant que foncteurs en $C_{.}$) sont encore exactement les foncteurs représentables sur $\mathrm{Ab}_{.}$.

[ \uline{NB}. Quand le foncteur $F$ a une structure plus riche, p.ex. une multiplication $F\otimes F\mathbin{\unsure{-|\to}}F$ i.e. est à valeurs dans les $\mathbb Z$-algèbres, \deleted{\unsure{est-ce|est $F($|et $F($}} cette structure sur $F(X_*)$ n'est pas reconstruite par
